\scanpage{491}
\mspage{249}

\uline{35} \uline{The injectivity of} $\Gamma_{\mathbb Q}\longrightarrow\mathrm{Autext}_{\mathrm{lac}}(\widehat{\Pi}_{0,3})$.

\margin{proved modulo lemma 2 below ---}

\uline{Theorem 1} Let $\Gamma=\mathrm{Gal}(\overline{\mathbb Q}_0/\mathbb Q)$, where $\overline{\mathbb Q}_0$ is the algebraic closure of $\mathbb Q$ in $\mathbb C$, let us consider the canonical homomorphism
\[
(1)\qquad \unsure{\Gamma}\xrightarrow{\Theta}\widehat{\widehat{\unsure{\mathcal T}}}_{0,3}=\mathrm{Autext}_{\mathrm{lac}}(\widehat{\Pi}_{0,3}).
\]
This homomorphism is \uline{injective}.

Proof.

\uline{Lemma 1} Let $x\in\mathbb Q$, $x\ne0,1$ i.e. $x\in U_{0,3}(\mathbb Q)$, let us choose a \deleted{loop} path on $\mathbb P^1(\mathbb C)$ from $P=-\unsure{\overline{j}|\overline{J}}$ (base point for defining $\Pi_{03}=\pi_1(U_{0,3}(\unsure{\mathbb C}),P)$ --- N.B. $U_{03}=\mathbb P^1_{\mathbb Q}\setminus\{0,1,\infty\}$) hence an isomorphism $\pi_1(U_{0,3\,\overline{\mathbb Q}_0},x)\simeq\widehat{\Pi}_{0,3}$, and by transport of structure an \inserted{actual} action of $\Gamma$ on $\widehat{\Pi}_{0,3}$ (not just outer). Then $K\unsure{=}\mathrm{Ker}\,\Theta$ operates trivially on $\widehat{\Pi}_{0,3}$.

\inserted{Proof. It suffices to see that the \uline{outer} operation of $K$ on $\pi_1(\overline V)$, where $V=U_{0,3}\setminus\{x\}$, is trivial.}
By the result of Белый, there exists a morphism (\uline{defined over} $\mathbb Q$, this is essential)
\[
\mathbb P^1_{\mathbb Q}\longrightarrow\mathbb P^1_{\mathbb Q},
\]
\deleted{\illegible{1}}\unsure{étale} over $U_{0,3}=\mathbb P^1_{\mathbb Q}\setminus\{0,1,\infty\}$, and such that $x\longmapsto0$. Let $U'$ be the induced étale covering of $U_{03}=U$, \deleted{\illegible{2}} $\pi'=\pi_1(\overline{U'})$ its geo-
\margin{490}

\scanpage{492}

metric fundamental group, on which $\Gamma$ hence $K\subset\Gamma$ operates by outer automorphisms. Then $\pi(\overline V)$ is a quotient of $\pi'$, respecting the outer operations of $\Gamma$, and it suffices to prove that $K$ operates trivially on $\pi'$. But $\pi'$ is identified with an open subgroup of $\widehat{\Pi}_{0,3}=\pi$ (respecting the outer operations of $K$, on which the outer operation of $K$ is trivial. It follows that the \inserted{outer} operation of $K$ on $\pi'$ factors through a \uline{finite} quotient group.

\margin{\diagram{\begin{tikzcd}[ampersand replacement=\&]
1 \arrow[r] \& \pi' \arrow[r] \arrow[d] \& \mathcal E'_K \arrow[r] \arrow[d] \& K \arrow[r,dash] \arrow[d,equal] \& 1 \\
1 \arrow[r,dash] \& \pi \arrow[r] \& \mathcal E_K \arrow[r,dash] \& K \arrow[r,dash] \& 1
\end{tikzcd}}}

[Indeed, the outer action of $K$ factors through the actual action of the extension group $\mathcal E_K$, which by the construction of $K$ is given by a homomorphism $\mathcal E_K\longrightarrow\pi/\unsure{Z|\mathcal Z}$ ($Z=\mathrm{Centre}(\pi)$), \deleted{\illegible{1}} the \unsure{induced|induced} action of $\mathcal E'_K$ is given by the composite $\mathcal E'_K\longrightarrow\mathcal E_K\longrightarrow\pi/Z$, whose image lies in $N/Z$, where $N$ is the normalizer of $\pi'$ in $\pi$. Thus the image of $K$ in $\mathrm{Autext}(\pi')$ is contained in that of $N/\pi'\longrightarrow\mathrm{Autext}(\pi')$, and $N/\pi'$ is finite ]
\margin{491}

\scanpage{493}
\mspage{250}

Returning to $\pi_1(\overline V)$, we thus see that the image of $K$ in $\mathrm{Autext}\,\pi_1(\overline V)$ is finite, hence the image of $K$ in $\mathrm{Aut}\,\widehat{\Pi}_{03}$ is finite. It moreover consists of inner automorphisms, so Lemma 1 will be a consequence of

\uline{Lemma 2} Every inner automorphism of finite order of $\widehat{\Pi}_{0,3}$ (free profinite group on two generators) is trivial. More precisely: $\widehat{\Pi}_{0,3}$ has a center reduced to $\{1\}$, and every element of $\widehat{\Pi}_{0,3}$ of finite order is reduced to 1.

This is a statement of pure profinite algebra, which I put off until later, to be done with the ``geometric'' part of the proof of the theorem.

\uline{Lemma} 3 For \unsure{every|every} \inserted{nonempty} open $\deleted{U}V\subset\deleted{\unsure{U_{0,3}}}\inserted{\mathbb P^1_{\mathbb Q}}$, \inserted{considering} the outer action of $\Gamma$ on $\pi_1(\overline{\deleted{U}V})$, its restriction to $K$ is trivial.

Passing if need be to a smaller
\margin{492}

\scanpage{494}

open, we may suppose by Белый that $\deleted{U}V$ is an étale covering of $U_{03}$ --- and the preceding argument then shows that the outer action of $K$ factors through a finite \inserted{quotient} group --- but this is not sufficient for our purpose, and does not by itself imply that this action is trivial.

\unsure{(}\deleted{But choose}\inserted{Besides, at the point I have reached,} we could have replaced $U$ by any quasi-projective smooth geometrically connected algebraic curve --- to find that $K$ operates on the outer group $\pi_1(\overline{\deleted{U}V})$ via a finite group. But here the hypothesis $\deleted{U}V\subset\mathbb P^1_{\mathbb Q}$ implies that there exists $y\in V(\mathbb Q)$, let $x\in U(\mathbb Q)$ be its image in $U=U_{0,3}$.

Taking \deleted{\unsure{the}} $y,x$ as base points for the fundamental groups, we now find an actual homomorphism of fundamental groups
\margin{493}

\scanpage{495}
\mspage{251}

\[
\pi'=\pi_1(\overline V,y)\hookrightarrow
\pi\;(\simeq\widehat{\Pi}_{0,3})=\pi_1(\overline U,x),
\]
compatible with an \uline{actual} action of $\Gamma$ on these groups. By lemma 1
\deleted{this}\inserted{the} induced action of $K$ on $\pi=\pi_1(\overline U,x)$ is trivial,
hence also its action on the subgroup $\pi'$. A fortiori the outer action is trivial, q.e.d.

We can now prove

\uline{Lemma 4} $K=(\unsure{1})$ (i.e. the theorem!)

Indeed, \inserted{under the conditions of lemma 2,} the action of
\deleted{\unsure{$\Gamma$}}\inserted{$K$} on the set
$I=S(\overline{\mathbb Q}_0)$, where $S=\mathbb P^1_{\mathbb Q}\setminus V$,
is deduced from the outer action of
\deleted{\unsure{$\Gamma$}}\inserted{$K$} on $\pi_1(\overline V)$,
\inserted{if $\operatorname{card} I\geq 2$,} as the action on the conjugacy classes of
``loop \unsure{subgroups|subgroups}''. As $K$ operates trivially on the outer \unsure{group},
it operates trivially on $I$. This \unsure{being} true for every $V$, we see that the action of $K$
on $\mathbb P^1(\overline{\mathbb Q}_0)$
\inserted{$=\overline{\mathbb Q}_0\cup\{\infty\}$}
is trivial, i.e. its action on $\overline{\mathbb Q}_0$ is trivial,
hence $K=(1)$, q.e.d.

\margin{Whew!}
\margin{494}

\scanpage{496}

It remains to set down the proof of lemma 2, which I am going to reformulate in a more general form:

\uline{Theorem 2} Let $\pi$ be a free profinite group on a set \inserted{finite $I$} of cardinality $\geq 2$.

Then:

a) The center of $\pi$ is \deleted{trivial} equal to $(1)$.

b) There is not \inserted{in $\pi$} any element of finite order other than $1$.

In the end I am stuck on a) -- just as I do not see why the center of $\Gamma$
(and of every open subgroup) must be reduced to $(1)$. Consult Deligne about this!

For b), here is an expedient.

Let $K$ be an algebraically closed field of characteristic $0$, and let us consider the function field
$L=K(x)$ of \deleted{\unsure{$U_{0,3\,K}$|$(U_{0,3})_K$}}
\inserted{$U$}$=\mathbb P^1_K\setminus$\deleted{\unsure{$I$|$\Sigma$}}
\inserted{($n+1$ points)}
\deleted{\unsure{($I$} \illegible{1} embedded in $\mathbb P^1(K)$)}
($n=\operatorname{card} I$). Then $\pi\simeq\pi_1(U)$, and it is a quotient of
$\mathcal E_L=\mathrm{Gal}(\overline L/L)$.
Since $\pi$ is free, $\mathcal E_L\longrightarrow\pi$ \unsure{lifts} to
$\pi\longrightarrow\mathcal E_L$, and it suffices to see

\margin{495}

\scanpage{497}
\mspage{252}

that $\mathcal E_L$ has no element of finite order $\ne 1$.

But according to Artin, there is no \unsure{automorphism} \inserted{\unsure{$\tau$}}
of finite order \inserted{$\ne 1$} of an algebraically closed field $\overline L$,
except in the case where $\overline L$ is a quadratic extension of a
\deleted{\illegible{1}} maximal \deleted{\unsure{real}}\inserted{ordered} field
\inserted{$R$}, \deleted{\unsure{fixed|finite}} which is the field of invariants of
$\tau$ (hence $\tau$ of order $2$ \unsure{exactly}).

\deleted{\illegible{1}} But in this \unsure{case|case} we would have $R\supset L$,
\inserted{and $L$}$\supset K$, now in $K$ the element $-1$ is a square,
so $R$ cannot be ordered; absurd!

\uline{Corollary 1 of \unsure{Theorem} 1} Let \unsure{$X$|$U$} be an algebraic curve
(smooth, \unsure{geometric\textsuperscript{ally}} connected, quasi-projective) over $k$ a finite extension of
$\mathbb Q$, \deleted{Then \unsure{there|A}} (I presume that the outer action of
$\mathcal E_k^{\overline k}$ on $\pi_1(\unsure{X|U}_{\overline k})$ is faithful
if $U$ \unsure{anabelian} -- failing to be able to prove it, I state:)
Then there exists a nonempty open subset $V$ of \unsure{$X$|$U$} such that
the \inserted{outer} action of $\mathcal E_k^{\overline k}$ on
$\pi_1(V_{\overline k})$ is faithful.

\margin{496}

\scanpage{498}

\uline{Proof}. By Белый, we know that we can find $V\subset\unsure{X}$, which is an étale covering of $U_{\deleted{\unsure{03}}}=(U_{03})_k$, I say that this $V$ will do. So let $K$ be the kernel of the outer \deleted{representation}\inserted{operation} of $\Gamma=\mathcal E_k^{\bar k}$ on $\pi'=\pi_1(V_{\bar k})$. Let $y$ be a closed point of $V$, rational over the finite extension $k'$ of $k$, corresponding to an open subgroup $\Gamma'\inserted{=\mathcal E_{k'}^{\bar k}}$ of $\Gamma$, let $K'=\Gamma'\cap K$. Let $x$ be the image of $y$ \unsure{in} $U$, we have
\[
\pi'=\pi_1(V_{\bar k},y)\hookrightarrow\pi=\pi_1(U_{\bar k},x),
\]
and by construction the outer action of $K'$ on $\pi'$ is trivial, hence $K'$ operates on $\pi'$ by inner automorphisms. Now we have

\margin{proved only modulo verification of $(2),(4)$ below -- for the \unsure{corollary} $(2)$ is moreover sufficient\ldots}
\uline{Lemma} 5 Let $\pi$ be a \inserted{profinite} group with loops, $\pi'$ an open subgroup, then every automorphism $u$ of $\pi$ which leaves $\pi'$ stable and induces on $\pi'$ an inner automorphism (resp. the identity) is inner (resp. the identity).

\margin{497}

\scanpage{499}
\mspage{253}

Let us admit this lemma for the moment -- it follows that the elements of $K'$, which operate on $\pi$ inducing inner automorphisms on $\pi'$, \unsure{induce|induce themselves} inner automorphisms on $\pi$ itself, i.e. that the outer action of $K'$ on $\pi$ is trivial. By theorem 1, we know on the other hand that the outer action of $\Gamma'$ on $\pi$ is faithful, hence $K'=(1)$. It follows that $K$ is \uline{finite}. But by Artin we know that the only finite subgroups \inserted{$\ne(1)$} of $\Gamma$ are those generated by a ``complex conjugation \deleted{of $\tau$}'' $\tau$, \unsure{corresponding|corresponding} \unsure{to a|to the} \deleted{real}\inserted{\unsure{ordered|ordered}} \unsure{maximal|maximal} subfield between $k$ and $\bar k$. But for such a $\tau$ we have $\chi(\tau)=-1$, hence the outer action of $\tau$ cannot be trivial -- we win. In fact, the proof has shown this:

\uline{Corollary 2} Let $\unsure{k|K}$ be \unsure{a} \deleted{extension of}\inserted{field of characteristic $0$},

\margin{498}

\scanpage{500}

\deleted{finite \unsure{type} over $\mathbb Q$}\inserted{$U$ a curve} algebraic (smooth etc.) over $k$, \unsure{such|such} that $\mathcal E_k^{\bar k}=\Gamma$ operates faithfully on $\pi_1(U_{\bar k})$ (outer operation).

Then for every geometric\textsuperscript{ally} connected étale covering $V$ of $U$, the outer operation of \deleted{\unsure{$\pi$}} $\Gamma$ on $\pi_1(V_{\bar k})$ is also faithful.

Remains to prove lemma 5. It suffices to prove it in the resp. case -- the other follows from it at once. \deleted{But let}\inserted{If $\pi$ has at least one class of loops (case of a curve} \deleted{$(x_i,y_i)_{1\le i\le g}$, $(l_j)_{1\le j\le\nu}$ a basis of $\pi$, ---}\inserted{algebraic affine i.e. nonproper), then $\pi$ is} \deleted{we shall admit that if} free -- and lemma 5 \deleted{\unsure{is found|is proved}}\inserted{is valid} precisely for such groups. If $(\unsure{\ell_i|h_i})_{1\le i\le\nu}$ is a system of generators, let $L_i$ be the closed subgroup generated by $\unsure{\ell_i|h_i}$, we shall admit that $\forall n\in\mathbb N^*$
\margin{not proved!}
\[
(2)\qquad L_i=\operatorname{Centr}_{\pi}(\unsure{L_i^n|\ell_i^n})
\]

This granted, if the automorphism $u$ of $\pi$ induces the identity on $\pi'$, $\exists n\in\mathbb N^*$ such that $\forall 1\le i\le\unsure{\nu}$
$u(\unsure{\ell_i|h_i}^{n})=\unsure{\ell_i|h_i}^{n}\unsure{{}_r}$, hence $u(\unsure{\ell_i|h_i})$ centralizes $\unsure{\ell_i|h_i}^{n}=u(\unsure{\ell_i|h_i})^{n}$,

\margin{499}